\section{Notation and conventions}\label{sec:notation}

\begin{table}[!htbp]
\centering
\caption{Principal symbols.}
\label{tab:symbols}
\small
\begin{tabular}{ll}
\toprule
Symbol & Meaning\\
\midrule
$i\in\{1,2\}$, $j$, $k$ & scan, echo, pathway (configuration order)\\
$\TR$, $t_{ijk}$ & repetition time; echo time after the RF pulse (ms)\\
$\alpha_i^{\rm nom}$, $\alpha_i=\Bone\alpha_i^{\rm nom}$ & nominal and actual flip angle of scan $i$\\
$\Bone$ & transmit-field scale (dimensionless)\\
$\Mzero$ & equilibrium magnetisation times receive sensitivity ($C\times\Mzero$)\\
$\Tone$, $\Ttwo$, $\Ttwos$ & relaxation times (ms); $\Ttwos=1/(1/\Ttwo+\Rtwop)$\\
$\Rtwop$ & reversible (Lorentzian) decay rate (1/ms), $\Rtwop\ge0$\\
$\dw$, $\phi_0$ & off-resonance (rad/ms), global receive phase (rad)\\
$E_1$, $E_2$, $\rho$ & $\E^{-\TR/\Tone}$, $\E^{-\TR/\Ttwo}$, $\TR/\Tone$\\
$F_k$, $Z_k$ & transverse and longitudinal configuration states\\
$a_k$ & steady-state amplitude of $F_k$ just after the pulse, per unit $\Mzero$\\
$u$, $v$ & the two invariants of \cref{thm:uv}\\
$S_{ijk}$ & noise-free voxel signal\\
$\sigma$ & noise standard deviation per real/imaginary component\\
$\theta$ & fitted parameters (log of positive ones, see below)\\
$\Fi$ & Fisher information matrix\\
\bottomrule
\end{tabular}
\end{table}

\paragraph{Transverse magnetisation and phases.} $M^+=M_x+iM_y$. An RF pulse of flip
$\alpha$ and phase $\varphi$ is the rotation $R_z(\varphi)R_x(\alpha)R_z(-\varphi)$; from
equilibrium it gives $M^+=-i\E^{i\varphi}\sin\alpha$. One repetition of the unbalanced gradient
multiplies $M^+$ by $\E^{i\psi}$, where $\psi$ is the intravoxel phase, and free precession by
$\dw$ for a time $t$ multiplies it by $\E^{i\dw t}$ (the same sense). The RF phase is constant
($\varphi=0$ unless stated).

\paragraph{Units.} Times in ms, rates in 1/ms, $\dw$ in rad/ms ($1$\,Hz $=2\pi\times10^{-3}$
rad/ms), angles in radians in formulas and degrees in text.

\paragraph{Parameters.} The fitted vector is
\begin{equation}
  \theta=(\log\Mzero,\ \log\Bone,\ \log\Tone,\ \log\Ttwo,\ \log\Ttwos,\ \dw,\ \phi_0),
\end{equation}
with the last two absent for magnitude data. Positive scale parameters are logged
(\cref{sec:logparam}); $\dw$ and $\phi_0$ are not. Bounds and spreads of log-parameters are
\emph{log-parameter standard deviations}: for a small value $s$ this is the relative
standard deviation; in general a spread $s$ corresponds to a one-sigma factor $\E^{\pm s}$.

\paragraph{Signal-to-noise ratio.} We write $\mathrm{SNR}(\Mzero)=\Mzero/\sigma$ with
$\Mzero=1$ for the reference tissue. The brightest MPME image of the reference tissue has
$|S|\approx0.088\,\Mzero$, so the image SNR is about $0.088\,\mathrm{SNR}(\Mzero)$.

\paragraph{Reference tissue and protocol.} Unless stated otherwise: $\Tone/\Ttwo/\Ttwos=
1500/70/60$\,ms (the tissue used for protocol optimisation in~\cite{cheng2019}), $\Bone=1$,
and the published in-vivo protocol (\cref{tab:protocol}).
